% 076_RSAtest_En_ch.tex -- chapter content Doc. 076
% Empirical tests of period-finding methods
% Complete rewrite August 2026 (R75)

\section{Empirical Tests of Period-Finding Methods}

\subsection*{Summary}

This document describes a test series on integer factorisation via period
finding. Classical methods (trial division, Pollard~$\rho$) are tested against
resonance-based search methods with different grid widths. The result places
the grid parameters correctly: they are numerical sampling quantities, not
physical resonance parameters.

\begin{keyresult}[Separate parameters]
The grid width of the search methods is denoted $\sigma$ here (values $1/42$,
$1/100$, $1/1000$, $1/100000$). It is \emph{not} the FFGFT geometry parameter
$\xi = 4/30000$. Earlier versions used the same symbol for both; the
separation is necessary because the quantities are unrelated.
\end{keyresult}

\section{Test Setup}

\subsection{Methods tested}

\begin{center}
\begin{tabular}{lll}
\toprule
Designation & Principle & Parameter \\
\midrule
\texttt{trial\_division} & Trial division up to $\sqrt{N}$ & --- \\
\texttt{pollard\_rho} & Cycle detection (Floyd) & --- \\
\texttt{period\_classic} & Sequential period search & $\sigma = 1/100000$ \\
\texttt{period\_universal} & Sequential period search & $\sigma = 1/100$ \\
\texttt{period\_medium} & Sequential period search & $\sigma = 1/1000$ \\
\texttt{period\_special} & Sequential period search & $\sigma = 1/42$ \\
\bottomrule
\end{tabular}
\end{center}

The grid width $\sigma$ determines the step size in which the search space is
traversed. It bears no relation to physical eigenfrequencies.

\subsection{Test categories}

\begin{enumerate}
  \item \textbf{Small semiprimes} ($N < 1000$): complete factorisation
    expected, runtime comparison.
  \item \textbf{Medium semiprimes} ($10^3 < N < 10^6$): scaling behaviour.
  \item \textbf{Twin-prime products} ($N = p(p+2)$): behaviour with
    structurally adjacent factors.
  \item \textbf{Large semiprimes} ($N > 10^6$): expected failure of period
    search.
\end{enumerate}

\section{Results}

\subsection{Small semiprimes}

All methods factorise completely. Trial division is fastest at these sizes;
period search carries no advantage here, since the effort of searching for $r$
exceeds that of direct division.

\subsection{Scaling behaviour}

\textbf{[K]} With growing $N$ the success rate of period search declines. The
reason is the number of candidates to be distinguished: the period $r$ divides
$\lambda(N) = \mathrm{lcm}(p-1,q-1)$ and can take values of order $N$. A
sequential search therefore requires $O(N)$ steps -- considerably worse than
trial division with $O(\sqrt{N})$.

\subsection{Influence of the grid width}

\textbf{[K]} Different $\sigma$ values yield different success rates. These
differences are \emph{sampling artefacts}: a finer grid finds more candidates
but costs more steps. There is no \enquote{optimal} $\sigma$ in a physical
sense -- only a trade-off between resolution and effort.

\begin{keyresult}[Why $\sigma$ generates no resonances]
The bicoherence test on the Riemann zeros (Doc.~316, August 2026) shows: genuine
phase coupling occurs exclusively at prime powers ($b^2 = 0.76$--$0.85$), not at
composite numbers ($b^2 = 0.018$--$0.025$). All $\sigma$ values used have
composite denominators ($42 = 2\cdot3\cdot7$, $100 = 2^2\cdot5^2$,
$1000 = 2^3\cdot5^3$). In the spectrum they form \emph{mixed points} without a
coupling line of their own. Their effect is purely numerical.
\end{keyresult}

\subsection{Twin-prime products}

For $N = p(p+2)$ the factors lie close together. Trial division benefits (a
find near $\sqrt{N}$), period search does not -- the period $r$ depends on
$\mathrm{lcm}(p-1,p+1)$ and shows no particular pattern. The
\enquote{twin-prime optimisation} via a special grid width described in earlier
versions is not reproducible.

\section{Placement}

\subsection{Comparison of methods}

\begin{center}
\begin{tabular}{llll}
\toprule
Method & Principle & Closure condition & Complexity \\
\midrule
Trial division & Direct division & Discrete arithmetic & $O(\sqrt{N})$ \\
Pollard $\rho$ & Cycle detection & Birthday paradox & $O(N^{1/4})$ \\
Period search & Spectral search & Rationalised grid & $O(N)$, Weyl limit \\
Shor's QFT & Quantum interference & Discrete phase & $O((\log N)^3)$ on QC \\
\bottomrule
\end{tabular}
\end{center}

\textbf{[K]} Classical period search is inferior to established methods. It has
no complexity advantage; the previously claimed agreement with Shor's
$O((\log N)^3)$ applies only to the quantum case and is not established for a
classical method.

\subsection{The structural limit}

\textbf{[K]} Beyond practical inferiority there is a limit in principle.
Following the Weyl obstruction (Doc.~316), the spectrum of any compact
resonance space grows as $N(\lambda) \sim \lambda^{d/2}$, while the arithmetic
counting density grows as $T\ln T$. A finite resonance space cannot carry this
density -- independently of grid width and computational precision.

\subsection{Primes as generators}

\textbf{[K]} What actually carries the resonance structure are the primes. The
Euler product treats each prime on its own, without cross terms; the dual length
spectrum contains only $k \ln p$. The harmonic series shows the same structure
physically: in the series $n f_0$, exactly the prime orders introduce new tonal
qualities, while $4 = 2\cdot2$ and $6 = 2\cdot3$ are combinations.

A method that treats primes as orthogonal axes works structurally correctly. A
method that treats a composite grid width as a resonance parameter does not.

\section{Balance}

\begin{center}
\begin{tabular}{ll}
\toprule
Statement & Status \\
\midrule
Period finding is a spectral problem & \textbf{[K]} \\
Primes carry the resonance structure & \textbf{[K]} measured by bicoherence \\
Preparation accounts for $\sim$95\,\% of gate cost & \textbf{[K]} (Doc.~176) \\
Sequential period search is $O(N)$ & \textbf{[K]} measured \\
$\sigma$ values are sampling grids & \textbf{[K]} \\
$\sigma$ values generate eigenfrequencies & \textbf{[X]} composite \\
Classical period search reaches $O((\log N)^3)$ & \textbf{[X]} not established \\
Twin-prime optimisation via grid width & \textbf{[X]} not reproducible \\
Scaling solvable by computational precision & \textbf{[X]} Weyl obstruction \\
\bottomrule
\end{tabular}
\end{center}

\textbf{The load-bearing sentence.} The test series documents the behaviour of
sequential period-search methods. It shows no advantage over established
methods, and the observed differences between grid widths are sampling
artefacts. The resonance description itself remains correct -- but its carriers
are the primes, not the grid.

\section{Relation to the Corpus}

The placement draws on Doc.~316 (August 2026): bicoherence test, Weyl
obstruction, and the role of primes as generators. The revision is declared
under R75 in Doc.~190.

\section{Verification}

Reproducible via \texttt{s1\_periodensuche.py} (correctness, scaling,
$\sigma$ factorisation, comparison of methods) and
\texttt{s2\_grenzen.py} (limits of both approaches). Standard library only,
target values as assertions, seed 20780458.
